\subsection{可对角化代数与平展代数}

定义1.——设 A 是 K 上的一个代数；如果存在一个整数 \(n \geqslant 0\) 使得 A 同构于乘积代数 \(\mathbf{K}^n\) ，则称 A 是可对角化的。如果由 A 通过标量扩张得到的 L 上的代数 \(\mathbf{A}_{(L)}\) 是可对角化的，则我们说 A 被 K 的一个扩张 L 所对角化。如果存在 K 的一个扩张将 A 对角化，则我们称 A 是平展的。

我们回顾乘积代数 \(K^{n}\) 是向量空间 \(K^{n}\) 配备由以下定义的乘积

\[
\left(x _ {1}, \dots , x _ {n}\right) \cdot \left(y _ {1}, \dots , y _ {n}\right) = \left(x _ {1} y _ {1}, \dots , x _ {n} y _ {n}\right).\tag{5}
\]

如果 \(\varepsilon_1, \ldots, \varepsilon_n\) 是 \(\mathbf{K}^n\) 的典范基，则我们有

\[
\varepsilon_ {i} ^ {2} = \varepsilon_ {i}, \quad \varepsilon_ {i} \varepsilon_ {j} = 0 \quad \text { if } \quad i \neq j\tag{6}
\]

且 \(1 = \varepsilon_{1} + \dots +\varepsilon_{n}\)

K上的每个平展代数都是交换的，且在K上具有有限次数。

命题1. --- 设A是域K上有限次数n的代数；则以下条件等价：

\begin{enumerate}
\def\labelenumi{\alph{enumi})}
\item
  该代数 \(\mathbf{A}\) 是可对角化的。
\item
  存在 A 的一组基 \((e_1, \ldots, e_n)\) ，使得 \(e_i^2 = e_i\) 且 \(e_i e_j = 0\) （对 \(i \neq j\)）。
\item
  A 到 K 的 K-代数同态生成向量 K-空间 A 的对偶。
\item
  每个A-模都是K上维数为1的子模的直和。
\end{enumerate}

\(a)\) 与 \(b)\) 的等价性由公式(6)可得；另一方面，这 \(n\) 个投影 \(\mathbf{K}^n \to \mathbf{K}\) 是代数同态，因此 \(a)\) 蕴含 \(c)\) 。如果 \(c)\) 成立，则 \(\mathbf{A}\) 到 \(\mathbf{K}\) 的代数同态构成 \(\mathbf{A}\) 的对偶的一组基（V，第27页，定理1），我们用 \(u_1, \ldots, u_n\) 表示它们；则 \(a \mapsto (u_i(a))\) 是 \(\mathbf{A}\) 到代数 \(\mathbf{K}^n\) 上的一个同构，从而得 \(a)\) 。这样我们已建立了条件 \(a), b)\) 与 \(c)\) 的等价性。

假设 \(b)\) 成立，并设 \(\mathbf{M}\) 为一个 A-模；则位似 \((e_i)_{\mathbf{M}}\) （比为 \(e_i\)）是 \(\mathbf{M}\) 的投影，且 \(\mathbf{M}\) 是诸 \(e_i\mathbf{M}\) 的直和，它们都是子 A-模。因此我们可以假设存在一个指标 \(i\) 使得 \((e_j)_{\mathbf{M}} = 0\) （对 \(j \neq i\)）。于是 \(\mathbf{M}\) 的每一个向量子空间都是子 A-模，从而得 \(d)\) 。

反之，假设 \(d\) 成立并考虑 A-模 \(\mathbf{A}_{s}\) 。于是存在一个基 \((f_{i})\) （向量 K-空间 \(\mathbf{A}\) 的），使得 \(\mathbf{A}f_{i} = \mathbf{K}f_{i}\) （对 \(i = 1, \ldots, n\)）。在必要时将每个 \(f_{i}\) 替换为适当的标量倍数后，我们可以假设 \(1 = f_{i} + \cdots + f_{n}\) 。若 \(i \neq j\) ，则 \(f_{i}f_{j}\) 属于 \(\mathbf{A}f_{i} \cap \mathbf{A}f_{j} = \mathbf{K}f_{i} \cap \mathbf{K}f_{j}\) ，因此它为零。现在 \(f_{i} = f_{i}f_{1} + \cdots + f_{i}f_{n} = f_{i}^{2}\) ，从而得 \(b\) ）。

推论。—— 设 L 为 K 的一个扩域，H 为从 A 到 L 的代数同态之集。我们有 Card \(H \leq [A:K]\) ，当且仅当 A 被 L 对角化时取等号。若 A 被 L 对角化，则 H 是向量 L-空间 \(\operatorname{Hom}_{\mathbf{K}}(\mathbf{A},\mathbf{L})\) 的一组基。

向量空间 \(\operatorname{Hom}_{\mathbf{K}}(\mathbf{A},\mathbf{L})\) 在 L 上的维数为 \([A:K]\) ，由公式(2)，且 H 是 \(\operatorname{Hom}_{\mathbf{K}}(\mathbf{A},\mathbf{L})\) 的一个自由子集，由定理1（V，第27页）。因此我们有 Card \(H \leqslant [A:K]\) ，当且仅当 H 是 \(\operatorname{Hom}_{\mathbf{K}}(\mathbf{A},\mathbf{L})\) 的一组基时取等号。存在一个向量 L-空间的同构，记为 \(\pi: \operatorname{Hom}_{\mathbf{K}}(\mathbf{A}, \mathbf{L}) \to \mathbf{A}_{(\mathbf{L})}^{*}\) ，其特征为 \(u(x) = (\pi u)(1 \otimes x)\) （对 \(x \in \mathbf{A}\)），且 \(\pi\) 将 \(\mathcal{H}\) 映到集合 \(\mathcal{H}_{\mathbf{L}}\) （即 \(\mathbf{A}_{(\mathbf{L})}\) 到 L 的 L-代数同态组成的集合）上。最后，命题1中 \(a)\) 与 \(c)\) 的等价性表明，L 上的代数 \(\mathbf{A}_{(\mathbf{L})}\) 可对角化当且仅当 \(\mathcal{H}_{\mathbf{L}}\) 生成 L 上的向量空间 \(\mathbf{A}_{(\mathbf{L})}^{*}\) 。这就完成了推论的证明。

命题 2. --- 设 A 是 K 上的代数，且 \(\Omega\) 是 K 的代数闭扩张。以下断言等价：

\begin{enumerate}
\def\labelenumi{\alph{enumi})}
\item
  该代数 \(\mathbf{A}\) 是平展的。
\item
  存在一个有限次扩张，它使 A 对角化。
\item
  该扩张 \(\Omega\) （\(\mathbf{K}\) 的）使 \(\mathbf{A}\) 对角化。
\end{enumerate}

假设A是平展的。设 n 为 A 在 K 上的次数，设 L 为 K 的一个使 A 对角化的扩张，并设 H 为从 A 到 L 的代数同态之集。根据命题1的推论，我们有 Card H = n。另一方面，对每个 \(u \in H\) ，我们有 \([u(A):K] \leqslant n\) 。由 V，第18页，定理2，由 H 中元素的像生成的 L 的子扩张 \(L'\) 在 K 上是有限次的。由于存在 A 到 \(L'\) 的 n 个不同的同态，扩张 \(L'\) 使 A 对角化，根据命题1的推论1。这表明 a) 蕴含 b)。

由于 K 的每个有限次扩张都同构于 \(\Omega\) 的一个子扩张（V，第20页，定理1），b) 蕴含 c)。最后，c) 显然蕴含 a)。

